The most directly related work is Pereira et al.\ \cite{10.1145/3136014.3136031}, who benchmarked 27 programming languages using RAPL and found that C and Rust are among the most energy-efficient, while Python consumes approximately 76 times more energy than C. A 2021 follow-up \cite{PEREIRA2021102609} extended this to examine energy, time, and memory jointly. Both studies report results in joules and use a subprocess-wrapping approach that includes process startup in the measurement window. Van Kempen et al.\ \cite{11334459} later argued that many observed language energy differences are partially explained by execution time differences alone, recommending that future studies separate the two effects. Our stdin handshake protocol and per-algorithm breakdown directly address this critique.

On the carbon accounting side, Lannelongue et al.\ \cite{lannelongue2021} proposed the Green Algorithms framework for estimating computational carbon footprints, and the Green Software Foundation published the Software Carbon Intensity (SCI) specification \cite{grnsoftware2023}, both of which mentions the gCO$_2$e metric used in this study. Strubell et al.\ \cite{strubell-etal-2019-energy} demonstrated the scale of AI training carbon costs, helping establish carbon footprint as a first-class concern in computational research.

For energy measurement, Khan et al.\ \cite{10.1145/3177754} validated Intel RAPL against external power meters with errors below 2\%, and Hähnel et al.\ \cite{10.1145/2425248.2425252} demonstrated that RAPL measurements for sub-millisecond workloads are dominated by sensor granularity rather than actual algorithm energy, directly motivating our focus on large input size results.